\documentclass[12pt,a4paper]{article}
\usepackage{amsmath,amsfonts,amssymb}
\usepackage{graphicx}
\usepackage{hyperref}
\usepackage{enumitem}
\usepackage{geometry}
\geometry{margin=2.5cm}
\emergencystretch=2em

\begin{document}

\begin{center}
    {\Large \textbf{Master Thesis Project in Quantum Information Science and Technology}}\\[0.3cm]
    {\large \textbf{Reinforcement Learning of Entanglement in Two-Electron Quantum Dots: Decoherence and its Effect on an Entangled Quantum Sensor}}\\[0.2cm]
    \vspace{0.3cm}
    {Duration: One Year (60 ECTS)}
\end{center}

\vspace{1cm}

\section*{1. Introduction}

Electrons trapped above the surface of superfluid helium form a clean
and highly controllable platform for quantum information processing
and quantum sensing: there is no host lattice, very little charge
noise and no nuclear spins, and single-electron confinement, control
and readout in microfabricated traps have been demonstrated
experimentally~[1--5].  Two electrons in an electrostatic double-well
potential are the helium analogue of a semiconductor double quantum
dot~[6,7].  Their low-energy states are a spin singlet and a spin
triplet with one electron in each well, separated by a tunable
exchange (singlet--triplet) splitting $\Delta_{ST}$.  The singlet is
insensitive to a uniform magnetic field but is rotated into the
triplet by a field \emph{difference} between the wells.  An entangled
electron pair is therefore a natural differential sensor, a concept we
have recently proposed for particle-physics applications within
CERN's DRD5 initiative on quantum sensors~[8].

Two questions decide whether such a sensor is useful.  The first is
how to shape the trap and the control pulses so that the pair is
prepared in a highly entangled state.  The trap is defined by several
electrode voltages, the dependence of the two-electron states on them
is strongly nonlinear, and the optimization target can only be
evaluated through a many-body calculation.  This is a setting where
reinforcement learning (RL) has become a standard tool in quantum
technology~[9--13], and a first RL optimization of the trap has been
carried out for the sensor proposal~[8].  The second question is how
long the entanglement survives.  Magnetic and electric noise couple
the electron pair to its environment, and in a recent
study~[14] we have shown how the correlated two-electron structure
determines the relaxation rates, the loss of spin entanglement and the
accuracy of reduced sensor models.

So far these two lines of work are separate: the RL optimization
assumed a closed system, and the open-system study used a fixed trap.
This thesis project combines them.  The aim is to use reinforcement
learning to find trap configurations and control protocols that
maximize the entanglement of a two-electron quantum dot, to study the
decoherence of the resulting entangled states, and to quantify what
this decoherence means for the performance of the two-electron sensor.

\section*{2. Physical System and Scientific Motivation}

The two electrons are described by the Hamiltonian (units
$\hbar^2/2m_e=1$)
\[
  \hat H_S=\sum_{i=1}^{2}\Bigl(-\nabla_i^2+v(\boldsymbol r_i)\Bigr)
  +\frac{\kappa}{\sqrt{|\boldsymbol r_1-\boldsymbol r_2|^2+\epsilon^2}},
  \qquad
  v(x,y)=k\Bigl[\bigl(\tfrac{x}{d}\bigr)^2-a^2\Bigr]^2+\delta\,x+\tfrac12k_yy^2,
\]
where the barrier height $k$, the well separation $a$ and the tilt
$\delta$ are the control parameters of a model trap; in a device they
are replaced by the electrode voltages~[5,8].  The problem is solved
on a discrete-variable grid and by configuration interaction (CI) in a
basis of left- and right-localized orbitals.  The six lowest states
are the singlet $|S\rangle$, the triplet $|T_{0,\pm}\rangle$ and two
charge-transfer singlets $|S'\rangle,|S''\rangle$ in which both
electrons occupy the same well.  A magnetic-field gradient
$\delta B$ enters as $H_B=h_\delta\,(S^z_L-S^z_R)$ with
$h_\delta=\tfrac12 g\mu_B\,\delta B$ and couples $|S\rangle$ to
$|T_0\rangle$ with the matrix element
$\mathcal M_{ge}=\langle T_0|S^z_L-S^z_R|S\rangle$.

The environment is included through a Lindblad master equation,
\[
  \frac{d\rho}{dt}=-i\bigl[H_S+H_B,\rho\bigr]
  +\sum_{\alpha,\omega}\gamma_\alpha(\omega)
  \Bigl(A_\alpha(\omega)\,\rho\,A_\alpha^\dagger(\omega)
  -\tfrac12\bigl\{A_\alpha^\dagger(\omega)A_\alpha(\omega),\rho\bigr\}\Bigr),
\]
where the coupling operators are the differential spin operator
$S^z_L-S^z_R$ (magnetic noise that differs between the wells) and the
charge imbalance $n_L-n_R$ (electric-field noise), and the rates
$\gamma_\alpha(\omega)$ follow from the spectral density of the bath
at the transition frequencies of the pair~[14,15].

Three quantities, all available from existing codes, connect the two
parts of the project.
\begin{itemize}
  \item \textbf{Accessible entanglement.}  A spin readout sees only
  the part of the state with one electron in each well.  Projecting
  the two-electron density matrix on this charge sector gives a
  two-spin state $\varrho$ with weight $p_{11}=\mathrm{Tr}\,\varrho$,
  and the entanglement figure of merit is
  $\mathcal E(\rho)=p_{11}\,C(\varrho/p_{11})$, with $C$ the
  Wootters concurrence~[16,17].
  \item \textbf{Decoherence times.}  The relaxation time
  $T_1=1/\Gamma$ with
  $\Gamma=\gamma(\Delta_{ST})\,\mathcal M_{ge}^2$, the dephasing time
  $T_2^*$ caused by slow fluctuations of the trap parameters, and the
  first disentanglement time $t_\times$ at which the concurrence of an
  initially entangled pair vanishes (entanglement sudden
  death~[18]).  In the reference trap of Ref.~[14]
  ($\Delta_{ST}=h\times 11.4$~MHz, $p_{11}^S=0.79$,
  $\mathcal M_{ge}=0.88$) the charge admixture of the singlet delays
  $t_\times$ by $16\%$ relative to ideal spins.
  \item \textbf{Sensitivity.}  The quantum Fisher information (QFI)
  $F_Q(t)$ of the state with respect to $h_\delta$ bounds the
  precision of any measurement through the quantum Cram\'er--Rao
  bound~[19--21].  In the noiseless limit and for vanishing exchange
  the entangled pair reaches $F_Q=4\mathcal M_{ge}^2t^2$, twice the
  value $2t^2$ of two independent spins.  A finite $\Delta_{ST}$ and
  decoherence both limit $F_Q$, which then has a maximum at a finite
  interrogation time.
\end{itemize}

These quantities pull in different directions.  A low barrier gives a
large exchange splitting, which makes the singlet easy to prepare and
the entangling exchange gate fast, but it also increases the charge
admixture (smaller $p_{11}$ and $\mathcal M_{ge}$), the sensitivity to
electric noise and to fluctuations of the barrier, and it suppresses
the sensing signal.  A high barrier gives nearly ideal spins and an
optimal signal, but an exponentially small exchange splitting, so that
the singlet can no longer be prepared by cooling.  It is well known
from quantum metrology that the most entangled state is not
necessarily the best probe once noise is included~[22--24].  The
project is organized around three questions:
\begin{enumerate}
  \item Which trap configurations and time-dependent control protocols
  maximize the accessible entanglement $\mathcal E$, and how
  efficiently does an RL agent find them compared with standard
  optimization and optimal-control methods?
  \item Through which channels, and how fast, do these optimized
  entangled states decohere?
  \item Does the configuration that maximizes the entanglement also
  maximize the sensitivity of the sensor when decoherence,
  preparation and readout are taken into account, and when does the
  entangled pair outperform two independent spins subject to the same
  noise?
\end{enumerate}

\section*{3. Objectives of the Thesis}

\subsection*{Objective 1: Two-electron structure and open-system dynamics}
\begin{itemize}
  \item Study the theory of electrons on helium, double quantum dots,
  CI methods, the Lindblad equation, entanglement measures and quantum
  parameter estimation.
  \item Work through the existing Python codes and Jupyter notebooks
  (sine-DVR, localized orbitals, spin-adapted CI, finite-bath exact
  dynamics, Lindblad solver) and reproduce the reference results of
  Ref.~[14]: the low-energy spectrum, $\Delta_{ST}$, $p_{11}$,
  $\mathcal M_{ge}$, the relaxation of $|T_0\rangle$ and the
  disentanglement time $t_\times$.
\end{itemize}

\subsection*{Objective 2: Reinforcement learning of the trap (static design)}
\begin{itemize}
  \item Formulate the trap design as an RL environment: the action is
  the vector of trap parameters (or electrode voltages), the reward is
  the accessible entanglement of the prepared state,
  \[
     R(\boldsymbol a)=\mathcal E\bigl[\rho_{\rm prep}(\boldsymbol a)\bigr]
     -w_b\,p_b(\boldsymbol a)-w_\mu\,p_\mu(\boldsymbol a),
  \]
  with penalties $p_b$, $p_\mu$ for the loss of the double-well shape
  and for delocalized orbitals as in Ref.~[8].  For
  $\rho_{\rm prep}$ use the thermal state at the electron temperature
  $T$, for which the competition between $\Delta_{ST}/k_BT$ and
  $p_{11}$ gives a well-defined optimum.
  \item Implement an actor--critic agent with Proximal Policy
  Optimization (PPO)~[10] and compare with gradient-free and
  Bayesian optimization, and with the analytical optimum of the
  two-site Hubbard model.
  \item Repeat the optimization of Ref.~[8] for the device potential
  with the corrected two-electron code.
\end{itemize}

\subsection*{Objective 3: Reinforcement learning of control pulses (dynamic preparation)}
\begin{itemize}
  \item Extend the environment to a sequential decision problem: at
  each time step the agent chooses piecewise-constant values of the
  barrier and the tilt, the state is propagated with the Lindblad
  equation, and the reward is the accessible entanglement at the final
  time, with a penalty for population left in the charge-transfer
  states.
  \item Train agents to prepare the singlet and other entangled states
  from product states (exchange pulses, adiabatic passage from the
  doubly occupied well) in the presence of dissipation, and study how
  the learned pulses change with the noise strength.
  \item Benchmark against analytical exchange-gate pulses and
  gradient-based optimal control (GRAPE), and compare policy-gradient
  (PPO) and off-policy actor--critic agents~[25].
\end{itemize}

\subsection*{Objective 4: Decoherence of the entangled states}
\begin{itemize}
  \item Compute $T_1$, $T_2^*$, $t_\times$ and the time dependence of
  $\mathcal E$ for the optimized configurations, for magnetic and
  electric noise acting separately and together, and at finite
  temperature.
  \item Include slow fluctuations of the barrier and the tilt through
  the Hellmann--Feynman sensitivities
  $\partial\Delta_{ST}/\partial k$ and
  $\partial\Delta_{ST}/\partial\delta$, and identify sweet spots.
  \item Check the Lindblad description against numerically exact
  finite-bath dynamics in the parameter and time windows where the
  latter is converged.
  \item Map the trade-off between entanglement, charge admixture and
  coherence time over the trap parameters.
\end{itemize}

\subsection*{Objective 5: Effect on the two-electron quantum sensor}
\begin{itemize}
  \item Compute the QFI and the classical Fisher information of a
  singlet--triplet readout for the estimation of a field gradient
  under decoherence, and determine the optimal interrogation time.
  \item Use the information per cycle,
  $\max_t F_Q(t)/(t+t_0)$ with $t_0$ the preparation and readout time,
  as the figure of merit, including the success probability $p_{11}$
  of the charge projection.
  \item Compare the entangled probe with two independent spins
  (Ramsey protocol) subject to the same noise and the same resources.
  \item Train a noise-aware agent with this figure of merit as reward
  for a complete prepare--sense--read sequence (large exchange for
  preparation and readout, small exchange during the interrogation),
  and compare its optimum with the configurations of Objectives~2
  and~3 that maximize the entanglement.
\end{itemize}

\subsection*{Objective 6: Possible extensions}
\begin{itemize}
  \item Noise operators and noise spectra derived from the electrode
  geometry of a device; low-frequency ($1/f$) exchange noise.
  \item Non-Markovian environments beyond the Lindblad equation
  (chain mapping or TEMPO~[26]).
  \item Surrogate models (parametric matrix models~[27]) of the reward to
  speed up the RL training.
  \item More than two electrons and arrays of wells, in connection
  with ongoing multi-configuration time-dependent Hartree--Fock
  studies of four-well devices~[4].
\end{itemize}

\section*{4. Methods and Tools}

\begin{itemize}
  \item \textbf{Electronic structure and dynamics:} Python (NumPy,
  SciPy) codes and Jupyter notebooks developed for Refs.~[8,14],
  including a step-by-step tutorial notebook; performance-critical
  parts can be rewritten in C++ if needed.  For the time-dependent
  controls the Hamiltonian and the coupling operators are tabulated in
  a fixed low-energy basis as functions of the trap parameters.
  \item \textbf{Reinforcement learning:} PyTorch or JAX with a
  Gymnasium-type environment; the existing PPO implementation of
  Ref.~[8] (Julia) serves as a reference.
  \item \textbf{Useful background:} quantum mechanics, many-body
  physics (FYS4480), machine learning (FYS-STK4155) and quantum
  computing (FYS5419).
\end{itemize}
All codes and results will be documented in Jupyter notebooks so that
the thesis can be reproduced.

\section*{5. Work Plan and Milestones (One-Year Timeline)}

\subsection*{Semester 1 (Months 1--6)}

\textbf{Months 1--2: Theory and reproduction}
\begin{itemize}
    \item Literature study; tutorial notebook; reproduction of the
    reference spectrum and of the open-system dynamics of Ref.~[14].
\end{itemize}
\textbf{Deliverable:} Short report with the reproduced reference numbers.

\bigskip

\textbf{Months 3--4: Static RL design}
\begin{itemize}
    \item RL environment and PPO agent for the trap parameters;
    thermal-state entanglement landscape; comparison with standard
    optimizers and with the Hubbard model.
\end{itemize}
\textbf{Deliverable:} Working RL optimization of the model trap and
of the device potential.

\bigskip

\textbf{Months 5--6: Dynamic RL control}
\begin{itemize}
    \item Time-dependent environment with dissipation; preparation of
    entangled states; GRAPE benchmark.
\end{itemize}
\textbf{Deliverable:} Learned preparation pulses with and without
noise.

\subsection*{Semester 2 (Months 7--12)}

\textbf{Months 7--8: Decoherence study}
\begin{itemize}
    \item Decoherence times and entanglement dynamics of the optimized
    states; combined noise channels, finite temperature, slow
    parameter noise; validation against finite-bath calculations.
\end{itemize}
\textbf{Deliverable:} Trade-off maps of entanglement against coherence.

\bigskip

\textbf{Months 9--10: Sensor performance}
\begin{itemize}
    \item QFI and classical Fisher information under decoherence;
    information per cycle; comparison with independent spins;
    noise-aware RL optimization of the full sensing sequence.
\end{itemize}
\textbf{Deliverable:} Answer to the question of when the entangled
sensor has an advantage.

\bigskip

\textbf{Month 11: Extensions and consolidation}
\begin{itemize}
    \item One of the extensions of Objective~6, depending on progress;
    final production runs and figures.
\end{itemize}

\bigskip

\textbf{Month 12: Writing and final submission}
\begin{itemize}
    \item Completion of the thesis manuscript (methods and theory
    chapters are written continuously during the year) and
    preparation of the final presentation.
\end{itemize}
\textbf{Deliverable:} Completed thesis and defense presentation.

\section*{6. Expected Outcomes}

\begin{itemize}
    \item An open and documented RL framework for the design and
    control of two-electron quantum dots, with the many-body and
    open-system solvers as environment.
    \item Trap configurations and control pulses that maximize the
    accessible entanglement, with a quantitative comparison of RL and
    conventional optimization.
    \item A map of how entanglement, charge admixture and decoherence
    compete in a double-well trap.
    \item A noise-aware assessment of the two-electron sensor,
    including the conditions under which it outperforms independent
    spins.
    \item A master thesis and, depending on the results, a
    contribution to a scientific article.
\end{itemize}

\section*{7. References}

\begin{enumerate}
    \item P.~M.~Platzman and M.~I.~Dykman, \textit{Quantum Computing with Electrons Floating on Liquid Helium}, Science \textbf{284}, 1967 (1999); S.~A.~Lyon, \textit{Spin-based quantum computing using electrons on liquid helium}, Phys.\ Rev.\ A \textbf{74}, 052338 (2006).
    \item M.~I.~Dykman and J.~Pollanen, \textit{Electrons in Quantum Dots on Helium: From Charge Qubits to Synthetic Color Centers}, Entropy \textbf{27} (2025).
    \item N.~R.~Beysengulov \textit{et al.}, \textit{Coulomb Interaction-Driven Entanglement of Electrons on Helium}, PRX Quantum \textbf{5}, 030324 (2024).
    \item O.~Leinonen \textit{et al.}, \textit{Design and dynamics of two-qubit gates with motional states of electrons on helium}, Phys.\ Rev.\ A \textbf{113}, 032624 (2026); O.~Leinonen \textit{et al.}, \textit{MCTDHF for simulations of multi-qubit systems}, in preparation.
    \item K.~E.~Castoria \textit{et al.}, \textit{Sensing and Control of Single Trapped Electrons above 1 K}, Phys.\ Rev.\ X \textbf{15}, 041002 (2025); G.~Koolstra \textit{et al.}, \textit{Strong coupling of a microwave photon to an electron on helium}, Nat.\ Phys.\ (2026).
    \item J.~R.~Petta \textit{et al.}, \textit{Coherent Manipulation of Coupled Electron Spins in Semiconductor Quantum Dots}, Science \textbf{309}, 2180 (2005).
    \item G.~Burkard, T.~D.~Ladd, A.~Pan, J.~M.~Nichol and J.~R.~Petta, \textit{Semiconductor spin qubits}, Rev.\ Mod.\ Phys.\ \textbf{95}, 025003 (2023).
    \item M.~E.~Perruzza \textit{et al.}, \textit{Electrons on Helium and Entangled Quantum Sensors for Particle Physics}, in preparation (2026).
    \item R.~S.~Sutton and A.~G.~Barto, \textit{Reinforcement Learning: An Introduction}, 2nd ed., MIT Press (2018).
    \item J.~Schulman, F.~Wolski, P.~Dhariwal, A.~Radford and O.~Klimov, \textit{Proximal Policy Optimization Algorithms}, arXiv:1707.06347 (2017).
    \item M.~Bukov and F.~Marquardt, \textit{Reinforcement Learning for Quantum Technology}, arXiv:2601.18953 (2026).
    \item M.~Bukov \textit{et al.}, \textit{Reinforcement Learning in Different Phases of Quantum Control}, Phys.\ Rev.\ X \textbf{8}, 031086 (2018); M.~Y.~Niu, S.~Boixo, V.~N.~Smelyanskiy and H.~Neven, \textit{Universal quantum control through deep reinforcement learning}, npj Quantum Inf.\ \textbf{5}, 33 (2019).
    \item F.~Belliardo, F.~Zoratti, F.~Marquardt and V.~Giovannetti, \textit{Model-aware reinforcement learning for high-performance Bayesian experimental design in quantum metrology}, Quantum \textbf{8}, 1555 (2024).
    \item J.~D.~\'Alvarez-Cuartas, N.~R.~Beysengulov, M.~Hjorth-Jensen, G.~F.~Lange, F.~P.~Massel, J.~Pollanen and J.~H.~Reina, \textit{Controlled open-system descriptions of two interacting electrons in a double-well trap}, in preparation (2026).
    \item H.-P.~Breuer and F.~Petruccione, \textit{The Theory of Open Quantum Systems}, Oxford University Press (2002).
    \item W.~K.~Wootters, \textit{Entanglement of Formation of an Arbitrary State of Two Qubits}, Phys.\ Rev.\ Lett.\ \textbf{80}, 2245 (1998).
    \item F.~Troiani, A.~Secchi and S.~Pittalis, \textit{Spins extracted from fermionic states and their entanglement properties}, Phys.\ Rev.\ A \textbf{112}, 012808 (2025).
    \item T.~Yu and J.~H.~Eberly, \textit{Sudden Death of Entanglement}, Science \textbf{323}, 598 (2009).
    \item S.~L.~Braunstein and C.~M.~Caves, \textit{Statistical distance and the geometry of quantum states}, Phys.\ Rev.\ Lett.\ \textbf{72}, 3439 (1994).
    \item M.~G.~A.~Paris, \textit{Quantum estimation for quantum technology}, Int.\ J.\ Quantum Inf.\ \textbf{7}, 125 (2009).
    \item C.~L.~Degen, F.~Reinhard and P.~Cappellaro, \textit{Quantum sensing}, Rev.\ Mod.\ Phys.\ \textbf{89}, 035002 (2017).
    \item S.~F.~Huelga \textit{et al.}, \textit{Improvement of Frequency Standards with Quantum Entanglement}, Phys.\ Rev.\ Lett.\ \textbf{79}, 3865 (1997).
    \item B.~M.~Escher, R.~L.~de Matos Filho and L.~Davidovich, \textit{General framework for estimating the ultimate precision limit in noisy quantum-enhanced metrology}, Nat.\ Phys.\ \textbf{7}, 406 (2011).
    \item R.~Demkowicz-Dobrza\'nski, J.~Ko{\l}ody\'nski and M.~Gu\c{t}\u{a}, \textit{The elusive Heisenberg limit in quantum-enhanced metrology}, Nat.\ Commun.\ \textbf{3}, 1063 (2012).
    \item T.~Haarnoja, A.~Zhou, P.~Abbeel and S.~Levine, \textit{Soft Actor-Critic: Off-Policy Maximum Entropy Deep Reinforcement Learning with a Stochastic Actor}, Proc.\ ICML (2018).
    \item A.~Strathearn, P.~Kirton, D.~Kilda, J.~Keeling and B.~W.~Lovett, \textit{Efficient non-Markovian quantum dynamics using time-evolving matrix product operators}, Nat.\ Commun.\ \textbf{9}, 3322 (2018).
    \item P.~Cook, D.~Jammooa, M.~Hjorth-Jensen and D.~D.~Lee, \textit{Parametric Matrix Models}, Nature Communications \textbf{16}, 5929 (2025).
\end{enumerate}

\end{document}
